% !TeX spellcheck = en_EN-EnglishUnitedKingdom
\chapter{Discussion}
This thesis aimed to understand the impact of assembly strategy on SynCom behaviour. After reconstructing and curating genome-scale metabolic models, simulating on apoplastic barley root medium showed that the models captured core metabolic functions and underlined the importance of community interactions in enabling growth, addressing the first research question. Analyses of niche overlap and cross-feeding behaviour highlighted different priorities between the communities and answering the second research question. Drop-out simulations demonstrated different robustness in growth to the removal of certain community members and revealed strain 892 (\textit{Bosea sp.} F3-2) as a potential keystone species, addressing the third research question. Together, these findings establish a basis for the answer of the fourth and fifth research questions, which require more literature context. 

\section{Metabolic interactions\textit{ in vivo} vs \textit{in silico}}

Comparing the \textit{in silico }findings to experimental results \textit{in planta}/\textit{in vitro} contextualises and potentially validates them. The two experimental frameworks investigate complementary interaction types: \cite{mahdi_low_2026}'s experiments primarily targeted antagonistic interactions, whereas GEM analysis focused on cooperative interactions. Consequently, the reconstructed communities show overall much weaker antagonistic interactions than reported \textit{in vitro}. For example, the drop-out simulations showed very few strains whose absence lead to widespread growth increases in others, underrepresenting the strongly antagonistic strains observed experimentally. Mechanisms of microbial antagonism are largely attributed to the production and secretion of antibiotics, bacteriocins, and toxins (\cite{garcia-bayona_bacterial_2018}), but the accurate representation of the producing biosynthetic pathways in GEMs remains a challenge. Databases lack information on reactions associated with secondary metabolism and these pathways are typically poorly annotated (\cite{qiu_flux_2023}). The representation of antagonistically acting metabolites in GEMs poses another challenge. Inhibiting a reaction upon the presence of a certain compound requires more complex constraints. As a result, the analysis conducted \textit{in silico }focussed on cooperative interactions which complicates the comparison to existing experimental data of the framework. Competition-based modelling approaches such as generalised Lotka-Volterra or MacArthur Consumer-Resource models could be used to complement the cooperative analyses and aid in the understanding of antagonistic interactions (\cite{van_den_berg_ecological_2022}).
\\


\textit{Hv}SC1 distributes its metabolic interactions across a broad set of smaller, more evenly-sized exchanges, while \textit{Hv}SC2 concentrates its interactions into fewer, higher-magnitude dependencies. \textit{Hv}SC1’s broad, shallow exchange strategy is, as expected, more robust to perturbation, since the loss of any single interaction means the loss of only a small fraction of a strain's total metabolite requirements, whereas a narrow, high-magnitude strategy creates more critical dependencies whose loss cannot easily be compensated by the remaining network.
One important assumption to keep in mind is that in all analyses, both \textit{in vitro }and\textit{ in silico}, all community members are equally abundant. While this framework allows constant evenness across both communities, it limits the applicability of both communities’ behaviours in practice, as natural systems are characterised by uneven and widely spread abundances (\cite{dawson_small_2017}).

\section{Keystone species}
Keystone species are defined as highly connected taxa, which play a unique and central role in community structure and function (\cite{banerjee_keystone_2018}). Because their influence is disproportionate to their abundance, removing them can cause profound systemic shifts (\cite{carlstrom_synthetic_2019}). Based on this framework, strain 892 can be termed a keystone species: It is a central exchange partner to its community members, its removal strongly decreases community growth, and its addition increases community growth. This pattern mirrors reports in other SynComs where a single strain's removal disproportionately destabilises the network (\cite{krumbach_metabolic_2021, sun_metabolic_2023}).
\\


Strain 892 has been identified as \textit{Bosea sp.} F3-2, and is a member of the \textit{Boseaceae} family. First isolated from the tobacco rhizosphere, this strain has been presented as a promising biocontrol candidate against quorum-sensing pathogens (\cite{zhang_aidb_2019}). Members of the genus \textit{Bosea} are classified as chemolithoheterotrophic, which is reflected in strain 892's metabolic profile. As a chemolithotroph, it is able to obtain energy from inorganic compounds, which has been proven for thiosulfate but has not been previously documented for the GEM-predicted nitrogen compounds (\cite{das_oxidation_1996}). Most notable is its involvement in heterotrophic nitrification, the oxidation of ammonium via nitrite to nitrate. This is consistent with literature emphasising how the substrate availability in the rhizosphere enables this pathway (\cite{martikainen_heterotrophic_2022}), through which strain 892 acts as a detoxifier by removing toxic nitrite and provides a nitrogen source to its fellow community members. 
The role as a central exchange partner is further supported by Strain 892's exclusive secretion of two siderophores. These compounds function as costly public goods that are frequently shared and also shape growth behaviour by mediating iron acquisition across the wider community (\cite{kramer_bacterial_2020}). 
\\


Strain 892’s beneficial role in a community is underscored by the increase in community growth upon the addition of strain 892 to the trait-optimised community. It facilitates growth for a subset of strains, for example by providing nitrate, while others grow less. The increased spread of individual growth rates upon 892's drop-in is notable because it runs counter to the compressive tendency of the L2-norm regularisation discussed below, which pushes growth rates toward equality rather than apart. This suggests that the spread in growth rates reflects a true biological effect of 892's addition. The change in growth behaviour is not purely caused by 892’s influence on metabolic exchanges, as the most frequent receivers are not necessarily the strains with the most improved growth. This implies that 892 reshapes community dynamics via more complex mechanisms than can easily be explained by its direct participation in pairwise cross-feeding analyses.
\\


In the \textit{in vivo }analysis, strain 892 is a persistent member of the core community over a six-day time period, and shows inhibition of only one strain in inhibition assays (\cite{mahdi_low_2026}). As the authors did not investigate detailed mechanisms of microbial cooperation or conduct community-wide drop-out experiments, strain 892 did not stand out. In the computational analysis, strain 892 is really only highlighted as a potential keystone upon its removal.  Besides the fact that strain 892 exchanges metabolites with all \textit{Hv}SC1 community members, it does not stand out as a structural outlier within its community when comparing its exchanges in number of metabolites, flux magnitude or symmetry. A validation of its keystone role in \textit{in vitro}/\textit{in planta} drop-out and drop-in experiments, followed by the analysis of changes in community behaviour and function would underline this predicted result.  
While it is clear that 892 has a special role within the communities, the exact underlying working mechanisms remain unclear. Results obtained from GEMs can be used to guide further experiments, as they highlight potential interesting behaviours which can be looked into experimentally. This will improve our understanding of keystone species, which remains a central task in plant-microbe modelling, as such keystone strains can be used to enhance host performance and community diversity (\cite{krumbach_metabolic_2021}).

\section{Community stability}
Efficient, innovative and sustainable applications of SynComs in agriculture require stable communities. It is therefore crucial to understand the drivers of stability and also how stability can be achieved in comparatively unstable communities. There are various ways to define stability in ecology. As assessing time-resolved processes with the used GEMs is hard, this thesis investigates stability indicators, namely balanced communities’ exchanges, and a community’s response to exclusions/additions. These further potentially drive stability, aiding stable composition over time, and help to improve the mechanistic understanding of long-term robustness. Literature highlights metabolic exchanges as a key driver of stable microbial communities (\cite{zelezniak_metabolic_2015}), consequently factors enhancing cooperation can also improve stability (\cite{wang_narrow-spectrum_2025}). Symmetry in both competitive and cooperative interactions promotes stability (\cite{butler_stability_2018}). Furthermore, it has been shown that communities with keystone taxa maintained more stable compositions and species abundances (\cite{liu_keystone_2025}). Based on this framework, the stability indicators of the two SynComs can be compared. 
\\


Both communities heavily rely on cross-feeding, however, the structure of the exchange networks differ. \textit{Hv}SC1 exchanges more metabolites per strain pair, showing greater average connectedness and lower, more symmetric exchange fluxes, while \textit{Hv}SC2's exchanges are comparatively less diverse and heavily asymmetric. Following \cite{butler_stability_2018}, this symmetry indicates greater stability for \textit{Hv}SC1, a prediction supported by the drop-out experiments, where only 892's removal noticeably perturbed \textit{Hv}SC1, while \textit{Hv}SC2's growth fluctuated more across all drop-outs. This implies that greater connectivity as well as symmetric and widely distributed interactions buffer metabolic needs and allow for adjustments to new conditions such as the removal of community members. Linking this back to functional redundancy, this shows that not all community members are strictly required for community stability, if a strain’s function is redundant with remaining members (\cite{allison_resistance_2008}).
\\


In both SynComs all members are equally abundant, hence species evenness, which can be one factor in establishing stability (\cite{wittebolle_initial_2009}), is also held constant across both communities. Consequently, evenness cannot account for the observed difference in robustness. The differences of \textit{Hv}SC1 and \textit{Hv}SC2 must therefore stem from the communities' other structural properties, such as connectivity or exchange flux symmetry. 
Keystone strains have been shown experimentally to support stability by (at least) two mechanisms: structural centrality, where highly connected keystone taxa provide redundant pathways that buffer the network against member loss (\cite{liu_keystone_2025}), and specialised function, where a keystone performs a metabolic role, that can not be substituted by other community members (\cite{iram_deciphering_2025, xun_specialized_2021}). Strain 892’s metabolic interactions exert influence on stability mainly through the latter: its role in nitrogen cycling, which exemplarily stands out upon its removal, demonstrated that its functional contribution can not be substituted by any other strain. 
\\


The investigated perspective on stability differs between the computational and experimental framework: While \cite{mahdi_low_2026} focused on stability over time, my analyses provide deeper understanding of stability indicators that may also aid establishing stability in the first place. Combining both approaches shows that the predicted stabilities of the \textit{in silico} communities align with the \textit{in planta} experimental results (\cite{mahdi_low_2026}). \textit{In planta}, \textit{Hv}SC1 largely retained its initial composition over a longer time, showed a larger number as well as more cooperative interactions and exerted higher environmental robustness. This is mirrored in the reconstructed \textit{Hv}SC1’s behaviour: the community is more balanced in its exchange fluxes as well as less impacted by single strain drop-outs, a result further supported by literature showing that an established community is more difficult to perturb (\cite{carlstrom_synthetic_2019}). In contrast, \textit{Hv}SC2 showed higher antagonism and major changes in community composition over time\textit{ in planta}, which is again consistent with the reconstructed \textit{Hv}SC2’s asymmetric exchanges and decreased robustness in drop-out experiments. All together, the results underline that symmetric metabolic interactions and the presence of stabilising keystone taxa are indicators of stable communities.

\section{Co-occurrence vs. trait-based: Impact of assembly strategies on community behaviour}\label{sec:assembly_strat_discussion}
The two SynComs clearly show how assembly strategy impacts community behaviour. The core community, which was assembled based on co-occurrence, demonstrates that co-occurrence promotes balanced community behaviour and robustness to changing conditions. Further, its more balanced individual growth rates are consistent with expectations from co-existence theory: selection acts against large asymmetries which would exclude slower-growing members long-term, to ensure stability over time (\cite{chesson_mechanisms_2000}). In \textit{Hv}SC1 these selection processes may already have taken place, whereas the trait-optimised community has not undergone this development, explaining why \textit{Hv}SC2 shows a more unevenly distributed exchange network with larger asymmetries in its individual growth rates. This distinction links nicely to a broader pattern in the literature: communities assembled for functional complementarity alone, without regard for ecological compatibility, tend to be less stable than those built around natural co-occurrence (\cite{de_souza_microbiome_2020}).
\\


The impact of assembly strategy goes beyond stability. Essential to a community’s effectiveness in promoting plant productivity and protection is its colonisation strength, efficient metabolite usage and balanced cooperations within the community. Literature shows that native core microbiomes often outperform other communities in exactly these aspects, resulting in greater potential to promote plant growth and yield (\cite{de_souza_microbiome_2020, zhou_superiority_2024}). 
\\


These results suggest a practical design approach: basing the design of trait-prioritised SynComs on native core communities that are commonly found across soils helps to overcome limitations in establishing stable host-microbiome relationships and scalability caused by narrow host and environment adaptability, i.e., a community only being stable in one specific environment (\cite{zhou_superiority_2024}). Further this would enhance the colonisation activity of SynComs, as outcompetition with naturally occurring bacteria can be prevented (\cite{de_souza_microbiome_2020}). The selected traits for \textit{Hv}SC2 by \cite{mahdi_low_2026} focused on plant-growth promoting or pathogen-protecting traits. In the future, the performance and stability of community-host associations could be improved by screening for strains that are enriched in genes that aid colonisation, such as those involved in boosted carbon metabolism (\cite{de_souza_genome_2019}). Ultimately these approaches point towards a hybrid assembly strategy, which combines core and trait-optimised communities, to help to realise the full potential of SynComs in agriculture. 

\section{Assumptions in model reconstruction and simulation}

\subsection{Model reconstruction quality and remaining curation challenges}
As shown by the MEMOTE scores, the curation process clearly improved the quality of the models and thus their accuracy in the representation of biological reality. Missing annotations cause the main gap to a “perfect score”. Although these gaps affect model quality metrics, annotations are less critical to functional predictions than network consistency and connectedness. 
\\


Comparing predicted to experimentally determined growth rates informs about plausibility of the models’ growth behaviour. In a community modelling context, this comparison is hard, as the community growth rate is not one experimentally measurable quantity but calculated based on different underlying metrics. The calculation approach, for example summing abundance scaled growth rates, substantially affects the community growth rate. Further, individual growth rates within the community context are difficult to compare to monoculture measurements in laboratory conditions, since a community model has access to a much larger combined resource pool than any single organism would in isolation. For reference, the fastest growing reported bacterium grows at a rate of $2.81 \,\text{h}^{-1}$ in optimal laboratory conditions (\cite{lee_functional_2019, lieven_memote_2020}), rhizospheric bacteria are reported with growth rates around $0.2-0.3 \,\text{h}^{-1}$ (\cite{lopez_growth_2023}). The abundance-scaled predicted growth rates fall within this range, but the unscaled ("raw") growth rates exceed biologically plausible values. Unrealistically high growth rates are one of many artifacts arising from uncertainties in various steps in GEM reconstruction and analysis. Uncertainties in genome annotation arise, as the mapping of homologs may be uncertain, or the reconstruction and interpretation of GPRs erroneous. This necessitates gap-filling, which is inherently uncertain because the reactions added are generally not supported by genomic evidence. One example, directly influencing growth rates, are sparsely characterised transport reactions, which lead to many pseudo-reactions with incomplete substrate specificity (\cite{bernstein_addressing_2021}). This causes thermodynamically impossible energy-generating cycles (EGCs), which allow for the production of energy metabolites (e.g., ATP or NADPH) without nutrient consumption, thereby inflating growth of a model. This partly explains the increased raw individual growth rates. One approach to resolve EGCs involves a variant of the GlobalFit algorithm, which suggests a minimal number of reaction removals that eliminate all EGCs. The algorithm is very computationally expensive and hence only applicable to small models (\cite{fritzemeier_erroneous_2017}). Even though this algorithm is almost 10 years old, there are no viable fully automated alternatives for large-scale GEMs, leaving the resolution of EGCs a matter of manual curation.
\\


Formulating the biomass objective function is another source of uncertainty. Very few experiments quantify the exact composition of a bacterium's biomass. Like many other tools, CarveMe uses a model organism’s biomass composition as a template, despite the importance of organism- or even strain-specific composition of the biomass (\cite{xavier_integration_2017}). Failure of menaquinol synthesis in strain 895 shows the problems with this generalised template approach. Although I only corrected the biomass function for this one strain, as all others functioned as expected, a thorough re-evaluation of all other biomass objective functions would ensure true biological replicability. Determining the exact biomass composition calls for large amounts of metabolomics data. Generating these datasets requires high-throughput workflows, such as automated, AI-guided robotic laboratory platforms (Jensen et al., 2025). The resulting experimental data can then be integrated into software tools like BOFdat (\cite{lachance_bofdat_2019}), which algorithmically infer accurate biomass stoichiometry directly from multi-omics data.
\subsection{Medium and its influence}
The composition of the growth medium has a major influence on bacterial growth rates \textit{in vitro} and impacts the behaviour of the communities \textit{in silico}. Consequently, defining a physiologically plausible and well-characterised medium is a crucial prerequisite for constraining community simulations. The medium used in this thesis is based on metabolomics data collected seven days post inoculation from the barley root apoplast by \cite{saake_ergosterol-induced_2025}. The authors also provided data from eight days post inoculation, which shows a substantial metabolic shift compared to the data obtained seven days post inoculation. These fast metabolic changes highlight the need for higher-resolution time-series datasets, and suggest averaging strategies across broader experimental replicates to construct robust baseline media definitions. Furthermore, spatial differences in metabolite abundance along the root may be another influence on experimental data (\cite{loo_sugar_2024}). Another challenge arises in the mapping of \textit{in planta }detected metabolites to their counterparts in the GEMs, as not all metabolites are listed in the BiGG database. Consequently, as the model reconstruction is based on BiGG, these metabolites can not be represented within the models. Additionally, not every metabolite found in the medium is present in the model, which requires gap-filling as well as including detailed knowledge about correctly absent pathways. During my medium generation process I identified 15 metabolites detected \textit{in planta}, which were not found in the GEMs. Moreover, defining realistic component uptake bounds remains challenging, as not much quantitative kinetic and transport data is available. These limitations highlight the need for more comprehensive metabolomic profiling and demonstrate that all results must be evaluated in the context of their environment.

\section{Limitations and relevance of the optimisation method} 
The results obtained in this thesis must be carefully interpreted, especially under the background of the employed constraint-based approach, underlying model quality and preceding assumptions. Contrasting to experiments \textit{in vivo}, the simulations do not include replicates. Even though the solution space is large, every optimisation returns a single maximal growth rate for every model, which might not capture natural variability. Consequently, standard statistical evaluation and significance testing, as typically applied to experimental datasets, were not performed here. Approaches such as flux variability analysis and alternative optima testing allow for statistical comparisons and present an aspect future work could investigate but were beyond the scope of this thesis.

\subsection{Limitations of (p)FBA}
The FBA method has its own limitations. Central to biologically relevant results is the choice of the objective function of a model. While for the exponentially growing community scenarios evaluated in this thesis, biomass maximisation remains a physiologically relevant proxy, depending on the research question this might change and is not achievable or preferable in every environmental condition (\cite{raman_flux_2009, schuetz_systematic_2007}). Furthermore, classical FBA operates under steady-state assumptions without integrating metabolic kinetics or enzyme allocation constraints. Consequently, the framework cannot account for proteomic capacity limitations, which can lead to unrealistic flux predictions. One example are non-growing or dormant strains carrying unphysiologically high metabolic fluxes. Merging FBA with kinetic methods in a dynamic FBA (dFBA) and non-linear optimisation approaches can help to overcome these limitations by predicting dynamic model behaviour (\cite{de_oliveira_nonlinear_2023, mahadevan_dynamic_2002}).
\subsubsection{Cooperative trade-off function - relevance of optimisation method}
The behaviour of the community is strongly influenced by the optimisation strategy employed in MICOM (\cite{diener_micom_2020}). Different growth rates were observed for different optimisation strategies (Fig. \ref{fig:hvsc1individualgrowthmethods}, \ref{fig:hvsc2individualgrowthmethods}), with profound implications for the entire system. A notable example is the blocking of sugar exchanges. The detected irregular rates observed under MICOM’s \texttt{cooperative\_tradeoff()} can be traced to its numerical crossover procedure, which is triggered when the optimisation solver fails to reach optimal convergence (see Section \ref{sec:Construction and simulation of community models}). The crossover heuristic attempts to recover a sparse flux distribution, but numerical inaccuracies during this step can artificially inflate flux throughput. When the crossover function was explicitly bypassed, no discrepancies in the results were detected. 
Aside from the computational aspects, the optimisation strategy should also reflect biological reality. When predicting metabolism \textit{in silico}, the flux distribution is closest to proteomics and transcriptomics measurements, when a GEM is optimised for a minimal sum of fluxes (pFBA) (\cite{lewis_omic_2010}), which serves as a first proxy for the assumption, that an organism will only invest minimal enzyme amounts. Further, growth rates simulated under L2 regularisation correlate more with experimentally observed replication rates compared to non-regularised results (\cite{diener_micom_2020}). 
\newpage

Therefore, I defined my own optimisation function, which similarly to MICOM’s original \linebreak \texttt{cooperative\_tradeoff()} function minimises the difference between individual growth rates (L2 norm) and also the sum of fluxes (pFBA) under a certain trade-off of maximal community growth. In contrast to MICOM’s function, it bypasses the crossover strategy and accepts non-optimal, yet numeric results. Numerical results in this case are acceptable, because small numerical deviations in individual growth rates do not meaningfully alter the overall community growth being compared. 
\\


The employed optimisation strategy reduces the number of non-growing strains compared to a regular pFBA, however it does not solve the problem of non/very little growing strains still carrying fluxes. While this is not a problem for the regular \textit{Hv}SC1 and \textit{Hv}SC2 communities, it certainly influences community behaviour of the drop-in \textit{Hv}SC2 + 892, as here three strains are growing less than $\mu_i = 0.001 \,\text{h}^{-1}$. As for the drop-in community, no detailed analysis of the exchange fluxes was performed, this was not specifically corrected for. One approach to address this in the future would be to shrink a strain’s upper and lower bounds once its growth falls below a defined threshold.
\\


A critical step in all these trade-off methods is the choice of the most suited trade-off parameter $\alpha$, which controls whether the optimisation prioritises community or individual growth. In my simulations, this parameter is always one, therefore optimising for maximal community growth. Decreasing the trade-off parameter yields overall smaller and more similar growth rates for the community members (Fig. \ref{fig:hvsc1individualgrowthtradeoff}, \ref{fig:hvsc2individualgrowthtradeoff}). As the goal was to compare individual growth rates in a community context, it is essential to choose a trade-off value that enables this comparison.  
\\


Even when using the custom function, certain observed patterns can still be artefacts of the trade-off strategy. Under the L2 norm, strains with very high growth rates are pushed down, while strains with very low growth rates are pushed up. Because this compression effect is distributed across all community members, its magnitude depends on community size: removing a member redistributes the compression across fewer remaining strains, weakening the push-down effect on strong growers and similarly the push-up effect on weak growers. Therefore, the strongest and weakest growing strains in a community should show a consistent directional shift upon any member's removal, as a mathematical consequence of the regularisation, independent of any true metabolic release or competition. This closely matches the pattern observed for the five strains in \textit{Hv}SC1 that responded consistently across all drop-out experiments. To further distinguish genuine biological truths from methodological artifacts, one should compare the results of different community modelling tools such as PyCoMo (\cite{predl_pycomo_2024}) or ReMIND (\cite{oftadeh_constraint-based_2024}). 

\section{Future perspectives}
This thesis helps to understand the impact of assembly strategies on community behaviour. Analysing differences in metabolic requirements, community interactions and perturbation robustness provides a detailed characterisation of the two investigated barley SynComs and aids future SynCom design. Nevertheless, there are plenty more open questions and tasks.
\\

 
One aspect is the improvement of model quality. Even with the high quality of the presented models, further refinement is required to reduce remaining errors and inaccuracies. Accuracy would be improved by removing EGCs, customising biomass functions, as well as including more experimental data on the individual strains, their environment and their interactions would further improve model and prediction accuracy. The inclusion of temporal processes, such as the previously mentioned dFBA would not only improve individual model quality, but also enable the long-term assessment of coexistence stability.
\\

 
Generally, the results of this thesis require experimental validation because GEM-based predictions rely on simplified assumptions that may not represent microbial behaviour in a real, dynamic environment. Especially drop-out and drop-in experiments of the keystone strain 892 would tell, if the predictions made here are also applicable \textit{in planta}, which would further help to understand the role of 892 in pathogen protection. Furthermore, experiments quantifying the cooperation between community members i.e, via systematic pairwise co-culture assays, would present valuable validation of the GEMs’ predictions.
\\


To understand how these SynComs protect the plant from pathogen invasion and promote growth, simulating the whole system, including GEMs of fungal pathogens and the barley plant would be necessary. While this comes with its own challenges, such as ensuring compatible metabolite exchanges, integrating host and pathogen models with existing bacterial community models would provide important insights into plant protection and microbiota function. The obtained predictions could easily be validated using the established experimental system of \cite{mahdi_low_2026}. 
\\


Finally, the insights gained in this thesis are relevant for the efficient and successful design of SynComs. As described in Section \ref{sec:assembly_strat_discussion}, combining core with trait-selected strains improves community stability and overall function. Understanding how core- and trait-based communities behave in isolation, as investigated in this thesis, is a foundational step towards understanding how such combined approaches would work in practice. Further modelling efforts and machine learning can support SynCom assembly by aiding the analysis of high-dimensional microbial datasets, predicting microbiome stability, compatibility, and emergent functions of new combinations and supporting the identification of strains possessing target functions from large datasets (\cite{garces-ruiz_machine_2026}).

